% SPDX-License-Identifier: CC-BY-4.0
\documentclass[11pt,letterpaper]{article}

\usepackage[T1]{fontenc}
\usepackage{amsmath,amssymb}
\usepackage{tikz}
\usepackage{pgfplots}
\pgfplotsset{compat=1.18}
\usepackage{algorithm}
\usepackage{algpseudocode}
\usepackage{booktabs}
\usepackage{enumitem}
\usepackage[hidelinks]{hyperref}

\newcommand{\N}{\mathbb{N}}
\newcommand{\Z}{\mathbb{Z}}
\newcommand{\Q}{\mathbb{Q}}
\newcommand{\R}{\mathbb{R}}
\newcommand{\C}{\mathbb{C}}
\newcommand{\Map}[3]{#1 \colon #2 \to #3}
\newcommand{\Mod}[1]{#1\text{-}\mathrm{Mod}}
\newcommand{\GroupHom}{\mathrm{Hom}_{\mathbf{Grp}}}
\newcommand{\RingHom}{\mathrm{Hom}_{\mathbf{Ring}}}

\newtheorem{theorem}{Theorem}
\newtheorem{lemma}[theorem]{Lemma}
\newtheorem{corollary}[theorem]{Corollary}
\newtheorem{definition}[theorem]{Definition}
\newtheorem{remark}[theorem]{Remark}
\newenvironment{proof}{\noindent\textit{Proof.}}{\hfill\(\square\)}

\title{Generalized Collatz Library}
\author{}
\date{\today}


\begin{document}
\maketitle
\section*{An Equal-Halves Gap with Three Crossings}

\begin{abstract}
For two equal-length shortcut words with equal odd counts, move one odd
letter across at most three common odd letters. The difference of their
ordered numerators is never divisible by the equal-halves sum modulus.
The proof is uniform in the word length and in all exterior support
positions. It uses a size bound except in two endpoint configurations,
which are closed by congruences modulo four and eight.
\end{abstract}

\section{Ordered supports}

For a finite support \(S\subseteq\{0,\ldots,h-1\}\), define
\[
A(S)=\sum_{j\in S}3^{\#\{k\in S:k>j\}}2^j.
\]
This is the ordered additive numerator of the shortcut word with odd
letters exactly at \(S\). Let \(0\le a<b<h\), and let \(C\) be a common
support containing neither endpoint. Set
\[
U=C\cup\{a\},\qquad V=C\cup\{b\},\qquad
r=|C|+1,\qquad M=2^h+3^r.
\]
Let \(d=b-a\). Write the common interior positions as
\(a+\ell_1<\cdots<a+\ell_t\), so
\(1\le\ell_1<\cdots<\ell_t\le d-1\), and put
\(s=\#\{x\in C:x>b\}\).

% Unlike the zero-crossing expression, this retains every changed suffix count.
\begin{lemma}[Crossing factor]
One has
\[
A(V)-A(U)=2^a3^s\Theta,\qquad
\Theta=2^d+2\sum_{i=1}^t3^{t-i}2^{\ell_i}-3^t.
\]
Moreover, \(M\) is coprime to \(2^a3^s\).
\end{lemma}
\begin{proof}
The moved letter contributes \(2^a3^{s+t}\) in \(U\) and \(2^b3^s\)
in \(V\). Common positions below \(a\) or above \(b\) have unchanged
strict suffix counts, so their contributions cancel. The suffix count
of the \(i\)-th interior position is \(s+t-i\) in \(U\) and one larger
in \(V\). Its difference is therefore
\(2\cdot3^{s+t-i}2^{a+\ell_i}\). Factoring gives the formula.
Here \(h>0\), \(r>0\), so \(M\) is odd and nonzero modulo three.
It is consequently coprime to every product of powers of two and three.
\end{proof}

% The endpoint cases use the number of exterior slots, not an independent bound on r.
\begin{theorem}[Three-crossing exclusion]
If \(t\le3\), then \(M\nmid A(V)-A(U)\).
\end{theorem}
\begin{proof}
The crossing factor and coprimality reduce the obligation to
\(M\nmid\Theta\). In every case below, \(\Theta>0\).

For \(t=0\), one has \(\Theta=2^d-1<2^h<M\).

For \(t=1\), \(d\ge2\) and
\[
0<\Theta=2^d+2^{\ell_1+1}-3<2^{d+1}\le2^h<M.
\]

For \(t=2\), \(d\ge3\), \(\ell_1\le d-2\), and \(\ell_2\le d-1\).
Thus
\[
0<\Theta=2^d+6\cdot2^{\ell_1}+2\cdot2^{\ell_2}-9
\le\frac72\,2^d-9.
\]
If \(h-d\ge2\), this is smaller than \(2^h\).
Otherwise \(h=d+1\), so \(a=0\), \(b=h-1\), and there are no exterior
positions. Consequently \(r=3\) and \(M=2^{d+1}+27\).
The bound gives \(0<\Theta<2M\). Divisibility would force
\(\Theta=M\), equivalently
\[
3\cdot2^{\ell_1}+2^{\ell_2}=2^{d-1}+18.
\]
Reduction modulo four forces \(\ell_1=1\). Substitution then gives
\(2^{\ell_2}=2^{d-1}+12\), impossible because \(\ell_2\le d-1\).

For \(t=3\), \(d\ge4\) and
\[
\begin{aligned}
\Theta&=2^d+18\cdot2^{\ell_1}
              +6\cdot2^{\ell_2}+2\cdot2^{\ell_3}-27,\\
0<\Theta&\le\frac{23}{4}\,2^d-27.
\end{aligned}
\]
Indeed, the upper bound uses respectively
\(\ell_1\le d-3\), \(\ell_2\le d-2\), and \(\ell_3\le d-1\).
Positivity follows already at the smallest permitted positions:
\(2^d\ge16\), and the three positive terms are at least \(36,24,16\).
If \(h-d\ge3\), then \(\Theta<2^h<M\).

If \(h=d+1\), there are no exterior positions, so \(r=4\) and
\(M=2^{d+1}+81\). The bound gives \(\Theta<3M\). Both \(\Theta\) and
\(M\) are odd, so a positive integral quotient less than three can
only be one. Divisibility therefore requires \(\Theta=M\), or
\[
9\cdot2^{\ell_1}+3\cdot2^{\ell_2}+2^{\ell_3}=2^{d-1}+54.
\]

If \(h=d+2\), there is exactly one possible exterior support slot.
Thus \(r\) is four or five, and \(\Theta<2M\) again forces
\(\Theta=M\). For \(r=5\), this would give
\[
9\cdot2^{\ell_1}+3\cdot2^{\ell_2}+2^{\ell_3}
=3\cdot2^{d-1}+135,
\]
contradicting parity. For \(r=4\), the equation instead is
\[
9\cdot2^{\ell_1}+3\cdot2^{\ell_2}+2^{\ell_3}
=3\cdot2^{d-1}+54.
\]

Both remaining equations have the form
\[
9\cdot2^{\ell_1}+3\cdot2^{\ell_2}+2^{\ell_3}
=\eta\,2^{d-1}+54,\qquad \eta\in\{1,3\}.
\]
Since \(d\ge4\), reduction modulo four forces \(\ell_1=1\).
After subtracting eighteen, reduction modulo eight forces
\(\ell_2=2\). Subtracting another twelve gives
\[
2^{\ell_3}=\eta\,2^{d-1}+24>2^{d-1},
\]
contradicting \(\ell_3\le d-1\). Every case is excluded.
\end{proof}

\section{Recovery to cycle exclusion}

% The obstruction is used only for the explicitly specified one-letter displacement.
\begin{corollary}[Equal-half cycle exclusion]
No positive classical shortcut cycle has a concatenated word \(UV\)
whose equal-length halves have the supports above and \(t\le3\).
\end{corollary}
\begin{proof}
The two halves have length \(h\) and odd count \(r\). Their concatenated
numerator is \(3^rA(U)+2^hA(V)\).
A positive closed orbit requires its determinant
\(2^{2h}-3^{2r}\) to be positive and to divide that numerator.
In particular, \(M=2^h+3^r\) divides the numerator.
Reduction modulo \(M\), where \(2^h\equiv-3^r\), gives
\(M\mid3^r(A(U)-A(V))\). Coprimality to three then forces
\(M\mid A(U)-A(V)\), contrary to the theorem.
\end{proof}

The obstruction is a written proof. Within
\nolinkurl{QCD.NoCrossing.001/lean/}, the file
\nolinkurl{NoCrossing.lean} formalizes the zero-crossing case only.
For the two independent equal-halves congruences, see
\nolinkurl{EqualHalves.tex} in \nolinkurl{QCD.EqualHalves.001/tex/}.
No claim is made for four or more crossings or for several moved letters.

\end{document}
